\documentclass[aspectratio=169, 10pt]{beamer}

% ======================================================================
% PAQUETES Y TEMA NATIVOS (MÁXIMA COMPATIBILIDAD)
% ======================================================================
\usepackage[utf8]{inputenc}
\usepackage[spanish]{babel}
\usepackage{graphicx}
\usepackage{xcolor}
\usepackage{booktabs}
\usepackage{circuitikz}
\usepackage{tikz}
\usepackage{listings}

\usetheme{Boadilla}
\usecolortheme{seahorse}

\lstset{
    language=Python,
    basicstyle=\ttfamily\tiny,
    keywordstyle=\color{blue}\bfseries,
    commentstyle=\color{gray}\itshape,
    backgroundcolor=\color{gray!5},
    frame=single,
    breaklines=true
}

% ======================================================================
% METADATOS
% ======================================================================
\title[Bode y Twin-T Notch]{Semana 4 - Sesión 1\\Diagramas de Bode y Filtro Twin-T Notch}
\subtitle{Introducción a la Mecatrónica}
\author[M.Sc. Ing. B. Quiroga]{M.Sc. Ing. Bernardo Quiroga Turdera}
\institute[UCB Tarija]{Universidad Católica Boliviana ``San Pablo''}
\date{Semestre 2-2026}

\begin{document}

\begin{frame}
    \titlepage
\end{frame}

\begin{frame}{Estrategia de Reajuste y Evaluación (Hito 1 PFI)}
    \begin{table}[]
        \centering
        \small
        \begin{tabular}{lll}
            \toprule
            \textbf{Sesión} & \textbf{Contenido} & \textbf{Entregable / Hito} \\
            \midrule
            \textbf{Sem 4 - Lunes} & Teoría Bode en $j\omega$ y Lab Twin-T & Mediciones CSV y Script Bode \\
            \textbf{Sem 4 - Miércoles} & Cascada: Puente + Notch & Validación en banco \\
            \textbf{Sem 5 - Lunes} & \textbf{EVALUACIÓN SUMATIVA 1} & \textbf{Hito 1 Calificado} \\
            \bottomrule
        \end{tabular}
    \end{table}
    \vspace{0.3cm}
    \textbf{Estructura de Evaluación (100 pts):}
    \begin{itemize}
        \setlength{\itemsep}{2pt}
        \item \textbf{50\% Escrita:} Cálculo de $Z_{sensor}$ y trazado analítico de Bode.
        \item \textbf{50\% Práctica (Banco + Git):} Filtro operando en vivo, validación en Python e informe IEEE.
    \end{itemize}
\end{frame}

\begin{frame}{Ubicación en el PFI: El Front-End Pasivo}
    % REDISEÑO TIKZ CON COORDENADAS ESTRICTAS PARA EVITAR FALLOS
    \begin{center}
        \begin{tikzpicture}[auto, thick, 
            block/.style={draw, rectangle, minimum width=6cm, minimum height=0.8cm, align=center, fill=gray!10, font=\scriptsize}]
            
            \node[block] (cam) at (0,0) {\textbf{Termistor NTC} (Cámara Térmica)};
            \node[block] (puente) at (0,-1.6) {\textbf{PUENTE DE WHEATSTONE (Lab 1)}\\Portadora CA ($10\text{ kHz}$)};
            \node[block] (notch) at (0,-3.4) {\textbf{FILTRO TWIN-T NOTCH (Lab 2)}\\Elimina ruido de red ($50\text{ Hz}$)};
            
            \draw[->] (cam) -- node[right, align=left] {\tiny Rama 4 con parásitos} (puente);
            \draw[->] (puente) -- node[right, align=left] {\tiny Modulación ($10\text{ kHz}$)\\ \tiny + Ruido Acoplado ($50\text{ Hz}$)} (notch);
            \draw[->] (notch) -- ++(0,-1.2) node[below, align=center] {\textbf{[ SALIDA HITO 1 ]}\\ \tiny Señal limpia al In-Amp};
        \end{tikzpicture}
    \end{center}
\end{frame}

\begin{frame}{Fundamentación: Los Decibelios y las Décadas}
    \begin{columns}
        \begin{column}{0.45\textwidth}
            La ganancia fasorial se comprime en escala logarítmica para ver detalles críticos:
            $$|H(j\omega)|_{\text{dB}} = 20 \log_{10} |H(j\omega)|$$
            \vspace{0.2cm}
            \textbf{Escala de Frecuencia (Log):}
            \begin{itemize}
                \item \textbf{Década:} Factor de 10 (Ej. $1\text{ Hz}$ a $10\text{ Hz}$).
                \item La pendiente se lee en \textbf{dB/década}.
            \end{itemize}
        \end{column}
        \begin{column}{0.55\textwidth}
            \small
            \begin{table}[]
                \begin{tabular}{cc}
                    \toprule
                    \textbf{Relación Lineal} & \textbf{Magnitud dB} \\
                    \midrule
                    $1.000$ & \textbf{$0\text{ dB}$} (Paso directo) \\
                    $\frac{1}{\sqrt{2}} \approx 0.707$ & \textbf{$-3\text{ dB}$} (Corte de potencia) \\
                    $0.100$ & \textbf{$-20\text{ dB}$} (Queda 10\%) \\
                    $0.010$ & \textbf{$-40\text{ dB}$} (Queda 1\%) \\
                    $0.001$ & \textbf{$-60\text{ dB}$} (Notch severo) \\
                    \bottomrule
                \end{tabular}
            \end{table}
        \end{column}
    \end{columns}
\end{frame}

\begin{frame}{Deducción Matemática (Twin-T sin Laplace)}
    \begin{columns}
        \begin{column}{0.45\textwidth}
            % CIRCUITO COMPACTADO EN BEAMER PARA NO DESBORDAR LA COLUMNA
            \begin{center}
                \begin{circuitikz}[scale=0.65, transform shape]
                    \draw (0,1.5) to[sV, l=$V_{in}$] (0,0) node[ground]{};
                    \draw (0,1.5) to[short, -*] (1,1.5);
                    
                    \draw (1,1.5) -- (1,3.5);
                    \draw (1,1.5) -- (1,-0.5);
                    
                    \draw (1,3.5) to[R, l={$R$}] (3.5,3.5) coordinate(nA) to[R, l={$R$}] (6,3.5);
                    \draw (nA) to[C, l={$2C$}, *-] (3.5,1.5) node[ground]{};
                    \node[above=0.15cm] at (nA) {\textbf{Nodo A}};
                    
                    \draw (1,-0.5) to[C, l={$C$}] (3.5,-0.5) coordinate(nB) to[C, l={$C$}] (6,-0.5);
                    \draw (nB) to[R, l={$R/2$}, *-] (3.5,-2.5) node[ground]{};
                    \node[above=0.15cm] at (nB) {\textbf{Nodo B}};
                    
                    \draw (6,3.5) -- (6,1.5) -- (6,-0.5);
                    \draw (6,1.5) to[short, *-o] (7,1.5) node[right]{\large $V_{out}$};
                \end{circuitikz}
            \end{center}
        \end{column}
        \begin{column}{0.55\textwidth}
            Aplicando la LKC ($\omega_0 = \frac{1}{RC}$):
            \vspace{0.2cm}
            $$H(j\omega) = \frac{1 - \left(\frac{\omega}{\omega_0}\right)^2}{1 - \left(\frac{\omega}{\omega_0}\right)^2 + j 4 \left(\frac{\omega}{\omega_0}\right)}$$
            \vspace{0.2cm}
            \textbf{Comportamiento Físico:}
            \begin{itemize}
                \item \textbf{A $\omega_0$:} Interferencia destructiva total entre ramas ($\pm 90^\circ$). $H \to 0\text{V} (-\infty\text{ dB})$.
                \item \textbf{Efecto de Carga $R_L$:} El amortiguamiento sube a $4(1 + R/R_L)$. Atenuación cae drásticamente.
            \end{itemize}
        \end{column}
    \end{columns}
\end{frame}

\begin{frame}{Ejercicio Integrador Resuelto}
    \textbf{Datos:} $R = 14.3\text{ k}\Omega, C = 220\text{ nF}$. Ruido $1.8\text{ V}_{pico}$ a $50\text{ Hz}$, Portadora $150\text{ mV}_{pico}$ a $10\text{ kHz}$. Atenuación real medida $-46\text{ dB}$.
    \vspace{0.3cm}
    \begin{enumerate}
        \setlength{\itemsep}{8pt}
        \item \textbf{Sintonía:} $f_0 = \frac{1}{2\pi (14.3\text{k})(220\text{n})} = \mathbf{50.59\text{ Hz}}$.
        \item \textbf{Ruido Residual:} $\hat{V}_{ruido} = 1.8 \cdot 10^{-46/20} = \mathbf{9.02\text{ mV}_{pico}}$ (Menor que señal útil).
        \item \textbf{Transmisión Portadora:} $H(10\text{kHz}) = \frac{1 - 197.67^2}{1 - 197.67^2 + j4(197.67)} = \mathbf{0.999795}$. Salida = $\mathbf{149.97\text{ mV}_{pico}}$.
        \item \textbf{Carga $10\text{ k}\Omega$:} $\vert{}H_{cargado}\vert{} = \frac{0.0232}{9.606} = 2.415 \times 10^{-3} \implies \mathbf{-21.35\text{ dB}}$ (Falla sin buffer).
    \end{enumerate}
\end{frame}

\begin{frame}[fragile]{Instrucciones LTSpice y Extracción de Bode}
    \textbf{Configuración Clave en LTSpice:}
    \vspace{0.2cm}
    \begin{itemize}
        \setlength{\itemsep}{4pt}
        \item Clic en fuente de voltaje y poner en propiedades: \texttt{\color{purple}AC Amplitude = 1}
        \item Nombrar nodos con tecla \colorbox{gray!20}{\texttt{\footnotesize\textbf{n}}}.
        \item Agregar directiva de barrido con tecla \colorbox{gray!20}{\texttt{\footnotesize\textbf{.}}}: \texttt{\color{purple}.ac dec 1000 1 20k}
        \item Agregar parámetro de carga con tecla \colorbox{gray!20}{\texttt{\footnotesize\textbf{.}}}: \texttt{\color{purple}.step param RL\_val list 1Meg 10k}
    \end{itemize}
    \vspace{0.3cm}
    \textbf{Ejecución Práctica (Lab 2):}
    \begin{itemize}
        \item Armar en protoboard. Variar manualmente la frecuencia de $5\text{ Hz}$ a $10\text{ kHz}$.
        \item Completar tabla de valores, calculando $20\log_{10}(V_{out}/V_{in})$.
        \item Extraer CSV en osciloscopio y correr el Script de Python \texttt{lab2\_bode\_processing.py}.
    \end{itemize}
\end{frame}

\end{document}
